% problem_id: S001-lemma-resolution-property-one-sheaf
% source_id: S001 (Stacks Project)
% source_file: spaces-more-morphisms.tex
% statement_locator: spaces-more-morphisms.tex lines 13565--13571
% proof_locator: spaces-more-morphisms.tex lines 13573--13709
% env: lemma
% id_note: label 在本来源内唯一
% pairing: tight (verified)
% status: complete (statement + author proof verbatim + reason + locators)
%
\begin{lemma}
\label{lemma-resolution-property-one-sheaf}
In Situation \ref{situation-descend-resolution-property},
the algebraic space $X$ has the resolution property
if and only if $\pi_*\mathcal{I}$ is the quotient of a
finite locally free $\mathcal{O}_X$-module.
\end{lemma}

\begin{proof}
The pushforward $\pi_*\mathcal{G}$ of a finite type
quasi-coherent $\mathcal{O}_Y$-module $\mathcal{G}$ is a finite
type quasi-coherent $\mathcal{O}_X$-module by
Descent on Spaces, Lemma
\ref{spaces-descent-lemma-finite-over-finite-module}.
In particular, if $X$ has the resolution property, then $\pi_*\mathcal{I}$
is the quotient of a finite locally free $\mathcal{O}_X$-module by
Derived Categories of Spaces, Definition
\ref{spaces-perfect-definition-resolution-property}.

\medskip\noindent
Assume that we have a surjection $\mathcal{E} \to \pi_*\mathcal{I}$
for some finite locally free $\mathcal{O}_X$-module $\mathcal{E}$.
In the rest of the proof we show that $X$ has the resolution property
by reducing to the Noetherian case handled in
Lemma \ref{lemma-resolution-property-in-Noetherian-case}.
We suggest the reader skip the rest of the proof.

\medskip\noindent
A first reduction is that we may view $X$ as an algebraic space over
$\Spec(\mathbf{Z})$, see
Spaces, Definition \ref{spaces-definition-base-change}.
(This doesn't affect the conditions nor the conclusion of the lemma.)

\medskip\noindent
By Limits of Spaces, Lemma
\ref{spaces-limits-lemma-finite-in-finite-and-finite-presentation}
we can write $Y = \lim Y_i$ with $Y_i$ finite and of finite
presentation over $X$ and where the transition maps are closed
immersions. Consider the closed subspace $Z = V(\mathcal{I})$ of $Y$.
Since $\mathcal{I}$ is of finite type, the morphism $Z \to Y$ is of
finite presentation. Hence we can find an $i$ and a morphism
$Z_i \to Y_i$ of finite presentation whose base change to $Y$
is $Z \to Y$, see Limits of Spaces, Lemma
\ref{spaces-limits-lemma-descend-finite-presentation}.
For $i' \geq i$ denote $Z_{i'} = Z_i \times_{Y_i} Y_{i'}$.
After increasing $i$ we may assume $Z_i \to Y_i$ is a
closed immersion (of finite presentation), see Limits of Spaces, Lemma
\ref{spaces-limits-lemma-descend-closed-immersion}. Denote
$\mathcal{I}_i \subset \mathcal{O}_{Y_i}$ the ideal
sheaf of $Z_i$ and denote $\pi_i : Y_i \to X$ the structure morphism.
Similarly for $i' \geq i$. Since $Z = \lim_{i' \geq i} Z_{i'}$ we have
$$
\pi_*\mathcal{I} = \colim \pi_{i', *}\mathcal{I}_{i'}
$$
The transition maps in the system are all surjective
as follows from the surjectivity of the maps
$\pi_{i, *}\mathcal{O}_{Y_i} \to \pi_{i', *}\mathcal{O}_{Y_{i'}}$
and the fact that $Z_{i'} = Z_i \times_{Y_i} Y_{i'}$.
By Cohomology of Spaces, Lemma
\ref{spaces-cohomology-lemma-finite-presentation-quasi-compact-colimit}
for some $i' \geq i$ the map $\mathcal{E} \to \pi_*\mathcal{I}$
lifts to a map $\mathcal{E} \to \pi_{i', *}\mathcal{I}_{i'}$.
After increasing $i'$ this map
$\mathcal{E} \to \pi_{i', *}\mathcal{I}_{i'}$
becomes surjective (since if not the colimit of the cokernels,
having surjective transition maps, is nonzero).
This reduces us to the case discussed in the next paragraph.

\medskip\noindent
Assume $X$ is an algebraic space over $\mathbf{Z}$ and that
$Y \to X$ is of finite presentation. By absolute Noetherian
approximation we can write $X = \lim X_i$ as a directed limit,
where each $X_i$ is a quasi-separated algebraic space of finite type
over $\mathbf{Z}$ and the transition morphisms are affine, see
Limits of Spaces, Proposition \ref{spaces-limits-proposition-approximate}.
Since $\pi : Y \to X$ is of finite presentation we can find an $i$
and a morphism $\pi_i : Y_i \to X_i$ of finite presentation
whose base change to $X$ is $\pi$, see
Limits of Spaces, Lemma
\ref{spaces-limits-lemma-descend-finite-presentation}.
After increasing $i$ we may assume $\pi_i$ is finite, see
Limits of Spaces, Lemma
\ref{spaces-limits-lemma-descend-finite}.
Next, we may assume there exists a finite locally
free $\mathcal{O}_{X_i}$-module $\mathcal{E}_i$
whose pullback to $X$ is $\mathcal{E}$, see
Limits of Spaces, Lemma \ref{spaces-limits-lemma-descend-invertible-modules}.
We may also assume there is a map
$\mathcal{E}_i \to \pi_{i, *}\mathcal{O}_{Y_i}$
whose pullback to $X$ is the composition
$\mathcal{E} \to \pi_*\mathcal{I} \to \pi_*\mathcal{O}_Y$, see
Limits of Spaces, Lemma
\ref{spaces-limits-lemma-descend-modules-finite-presentation}.
The cokernel
$$
\mathcal{E}_i \to \pi_{i, *}\mathcal{O}_{Y_i} \to \mathcal{Q}_i \to 0
$$
is a coherent $\mathcal{O}_{Y_i}$-module whose pullback to $X$
is the (finitely presented) cokernel $\mathcal{Q}$ of the map
$\mathcal{E} \to \pi_*\mathcal{O}_Y$. In other words, we have
$\mathcal{Q} = \pi_*(\mathcal{O}_Y/\mathcal{I})$. Consider
the map
$$
\mathcal{E}_i
\otimes_{\mathcal{O}_{X_i}}
\pi_{i, *} \mathcal{O}_{Y_i}
\longrightarrow
\pi_{i, *}\mathcal{O}_{Y_i}
\otimes_{\mathcal{O}_{X_i}}
\pi_{i, *} \mathcal{O}_{Y_i}
\to
\pi_{i, *} \mathcal{O}_{Y_i}
\to
\mathcal{Q}_i
$$
where the second arrow is given by the algebra structure
on $\pi_{i, *}\mathcal{O}_{Y_i}$. The pullback of this map
to $Y$ is zero because the image of $\mathcal{E} \to \pi_*\mathcal{O}_Y$
is the ideal $\pi_*\mathcal{I}$. Hence by
Limits of Spaces, Lemma
\ref{spaces-limits-lemma-descend-modules-finite-presentation}
after increasing $i$ we may assume the displayed composition
is zero. This exactly means that the imag of
$\mathcal{E}_i \to \pi_{i, *}\mathcal{O}_{Y_i}$
is of the form $\pi_{i, *}\mathcal{I}_i$
for some coherent ideal sheaf $\mathcal{I}_i \subset \mathcal{O}_{Y_i}$.
Since $\mathcal{E}_i \to \pi_{i, *}\mathcal{O}_{Y_i}$
pulls back to $\mathcal{E} \to \pi_*\mathcal{O}_Y$ we
see that the pullback of $\mathcal{I}_i$ to $Y$ generates $\mathcal{I}$.
Denote $V_i \subset Y_i$ the open subspace whose complement
is $V(\mathcal{I}_i) \subset Y_i$. Then $V$ is the inverse image
of $V_i$ by the comments above. After increasing $i$ we may assume
that $V_i$ is affine and that $\pi_i|_{V_i} : V_i \to X_i$ is 
\'etale, see
Limits of Spaces, Lemmas
\ref{spaces-limits-lemma-limit-is-affine} and
\ref{spaces-limits-lemma-descend-etale}.
Having said all of this, we may apply
Lemma \ref{lemma-resolution-property-in-Noetherian-case}
to conclude that $X_i$ has the
resolution property. Since $X \to X_i$ is affine
we conclude that $X$ has the resolution property too
by Derived Categories of Spaces, Lemma
\ref{spaces-perfect-lemma-resolution-property-goes-up-affine}.
\end{proof}
